\documentclass[conference]{IEEEtran}
\usepackage{amsmath,amssymb,amsfonts}
\usepackage{algorithmic}
\usepackage{graphicx}
\usepackage{textcomp}
\usepackage{xcolor}
\usepackage{booktabs}
\usepackage{hyperref}
\usepackage{cleveref}

\def\BibTeX{{\rm B\kern-.05em{\sc i\kern-.025em b}\kern-.08em
    T\kern-.1667em\lower.7ex\hbox{E}\kern-.125emX}}

\begin{document}

\title{OrthoCache GPU: Fused Spectral KV-Cache Eviction\\with Split-K Parallelization on NVIDIA GPUs}

\author{
\IEEEauthorblockN{Justin Arndt}
\IEEEauthorblockA{justinarndt05@gmail.com \\
\url{https://github.com/j-arndt/orthocache-gpu}}
}

\maketitle

\begin{abstract}
We present OrthoCache GPU, a Triton-based implementation of spectral
KV-cache eviction for NVIDIA GPUs. OrthoCache uses the Fast Walsh--Hadamard
Transform (FWHT) to compute per-block spectral energy ratios entirely in
SRAM, skipping noise-dominated attention blocks before the quadratic
attention computation begins. Our GPU-native implementation introduces
two key contributions: (1)~a \textbf{fused God Kernel} that performs
spectral analysis, eviction scoring, and predicated attention in a
single Triton kernel launch with zero intermediate DRAM traffic, and
(2)~\textbf{Split-K parallelization} with interleaved (cyclic) tile
assignment that distributes the KV-cache workload across all streaming
multiprocessors. On an RTX 4060 Laptop GPU (24 SMs, Ada Lovelace),
the Split-K kernel achieves a \textbf{15.3$\times$ speedup} over
dense attention at 32,768 tokens, with \textbf{sub-linear latency
scaling} from 4K to 32K tokens (0.125\,ms $\to$ 0.172\,ms). The
algorithm's correctness is formally verified in Lean~4: Parseval's
identity for the WHT, the exponential Total Variation truncation
bound, and a hardware-native perfect eviction theorem under IEEE~754
float32 underflow semantics.
\end{abstract}

\begin{IEEEkeywords}
KV-cache eviction, attention, Walsh--Hadamard Transform, GPU, Triton,
spectral analysis, transformer inference, Split-K, formal verification
\end{IEEEkeywords}

% ===================================================================
\section{Introduction}\label{sec:intro}

Long-context large language model (LLM) inference is fundamentally
constrained by the memory bandwidth required to load the key-value
(KV) cache at each decoding step. At 32K tokens with 128-dimensional
heads, a single attention head's KV-cache occupies 16\,MB---and this
must be streamed from DRAM for every generated token.

OrthoCache~\cite{orthocache_tpu} addresses this by using the Walsh--Hadamard
Transform (WHT) to identify and skip semantically redundant KV-cache
blocks \textit{before} the quadratic attention computation. The original
implementation targets Google TPUs via Pallas/XLA. This paper presents
the GPU/Triton edition, which exploits NVIDIA's programmer-controlled
shared memory (SRAM) to achieve a fundamentally different---and
significantly more efficient---kernel architecture.

\subsection{Contributions}

\begin{enumerate}
\item \textbf{Fused God Kernel.} A single Triton kernel that performs
FWHT spectral analysis, $\zeta$-based eviction, and predicated
attention in one launch. Key tiles are loaded to SRAM once and reused
across both phases, eliminating the second DRAM load.

\item \textbf{Split-K Parallelization.} Distribution of KV-cache tiles
across all SMs with interleaved (cyclic) assignment and log-sum-exp
reduction, achieving near-perfect SM utilization.

\item \textbf{Sub-Linear Scaling.} Empirical demonstration that
OrthoCache latency stays nearly flat from 4K to 32K tokens on
consumer GPU hardware (38\% increase vs.\ 488\% for dense attention).

\item \textbf{Lean~4 Verification.} Formal proofs of Parseval's
identity, the exponential TV truncation bound, and a novel
hardware-native perfect eviction theorem shared with the TPU edition.
\end{enumerate}


% ===================================================================
\section{Background}\label{sec:background}

\subsection{KV-Cache and the Memory Wall}

% TODO: Expand with quantitative memory bandwidth analysis

\subsection{Walsh--Hadamard Transform}

The Walsh--Hadamard matrix $H_n$ of dimension $2^n \times 2^n$ is
defined by the Kronecker recurrence:
\begin{equation}
H_0 = [1], \quad H_{n+1} = H_n \otimes \begin{bmatrix} 1 & 1 \\ 1 & -1 \end{bmatrix}
\end{equation}

Parseval's identity (Lean~4 verified):
\begin{equation}\label{eq:parseval}
\|H_n \cdot x\|^2 = 2^n \cdot \|x\|^2
\end{equation}


% ===================================================================
\section{Method}\label{sec:method}

\subsection{Multi-Band Spectral Scoring}\label{sec:multiband}

We decompose the WHT spectrum into four frequency bands and define
the spectral ratio:
\begin{equation}
\zeta = \frac{E_{\text{high}}}{E_{\text{low}}} = \frac{\sum_{j=48}^{63} |\hat{k}_j|^2}{\sum_{j=0}^{15} |\hat{k}_j|^2}
\end{equation}

Blocks with $\zeta > \zeta_{\max}$ are classified as noise-dominated
and evicted from the attention computation.

\subsection{Truncation Bound}\label{sec:bound}

\begin{theorem}[OrthoCache Truncation Bound, Lean~4 Verified]
Given logits $z : \text{Fin}(N) \to \mathbb{R}$, a retained set $S$,
and an evicted set $S^c$ where all evicted logits satisfy $z_i < \beta$
and $z_{\max} = \max_{j \in S} z_j$:
\begin{equation}
\text{TV}(\alpha, \hat{\alpha}) \leq |S^c| \cdot \exp(\beta - z_{\max})
\end{equation}
\end{theorem}

\subsection{Fused Kernel Architecture}\label{sec:kernel}

% TODO: Detailed kernel pipeline diagram

\subsection{Split-K Parallelization}\label{sec:splitk}

The kernel grid is $(\text{num\_heads}, \text{num\_splits})$. SM $s$
processes tiles $\{s, s+K, s+2K, \ldots\}$ where $K = \text{num\_splits}$.

The log-sum-exp reduction merges partial states:
\begin{align}
m_{\text{new}} &= \max(m_1, m_2) \\
l_{\text{new}} &= l_1 e^{m_1 - m_{\text{new}}} + l_2 e^{m_2 - m_{\text{new}}} \\
\text{acc}_{\text{new}} &= \frac{l_1 e^{m_1 - m_{\text{new}}} \cdot \text{acc}_1 + l_2 e^{m_2 - m_{\text{new}}} \cdot \text{acc}_2}{l_{\text{new}}}
\end{align}

\subsubsection{Interleaved vs.\ Contiguous Assignment}

% TODO: SM utilization comparison figure


% ===================================================================
\section{Lean 4 Formalization}\label{sec:lean}

Three proof modules verify the algorithm's mathematical guarantees:

\begin{enumerate}
\item \texttt{ParsevalWHT.lean}: $H_n^T H_n = 2^n I$ and Parseval's identity.
\item \texttt{TruncationBound.lean}: $\text{TV}(\alpha, \hat{\alpha}) \leq |S^c| \cdot \exp(\beta - z_{\max})$.
\item \texttt{QuantizedTruncation.lean}: When $z_{\max} - \beta \geq 88.72$, the quantized exponential evaluates to exactly zero, yielding $\text{TV} = 0$.
\end{enumerate}

All proofs are algorithm-generic---they hold over $\mathbb{R}$ and
general matrices, with no hardware-specific assumptions. The IEEE~754
threshold (88.72) applies identically to NVIDIA GPU float32 and
Google TPU float32.


% ===================================================================
\section{Experimental Results}\label{sec:results}

\subsection{Hardware and Setup}

All benchmarks run on an NVIDIA GeForce RTX 4060 Laptop GPU (Ada
Lovelace, SM~8.9, 24 SMs, 8\,GB GDDR6) with CUDA~12.4, PyTorch~2.6,
and Triton~3.1.

\subsection{Latency Comparison}

\begin{table}[h]
\centering
\caption{Dense vs.\ OrthoCache Split-K Attention Latency}
\label{tab:latency}
\begin{tabular}{@{}cccc@{}}
\toprule
Context & Dense & OrthoCache & Speedup \\
\midrule
1,024  & 0.125\,ms & 0.075\,ms & 1.67$\times$ \\
4,096  & 0.448\,ms & 0.125\,ms & 3.59$\times$ \\
8,192  & 0.807\,ms & 0.173\,ms & 4.66$\times$ \\
16,384 & 1.348\,ms & 0.117\,ms & 11.5$\times$ \\
\textbf{32,768} & \textbf{2.635\,ms} & \textbf{0.172\,ms} & \textbf{15.3$\times$} \\
\bottomrule
\end{tabular}
\end{table}

\subsection{Sub-Linear Scaling}

% TODO: Figure showing latency curves (dense linear vs ortho flat)

\subsection{SRAM Budget}

% TODO: Table showing per-buffer allocation

\subsection{DRAM Traffic Reduction}

% TODO: Traffic comparison table


% ===================================================================
\section{Cost-Benefit Analysis}\label{sec:cost}

% TODO: Summary from docs/cost_benefit_analysis.md


% ===================================================================
\section{Related Work}\label{sec:related}

\textbf{FlashAttention}~\cite{dao2023flashattention2} tiles attention
for IO efficiency but computes all blocks---no eviction.
\textbf{PagedAttention}~\cite{kwon2023pagedattention} manages KV-cache
memory via virtual paging but does not reduce compute.
\textbf{H$_2$O}~\cite{zhang2023h2o} uses heavy-hitter scoring for
eviction but lacks formal guarantees and requires a full attention
pass for scoring (circular dependency).
\textbf{StreamingLLM}~\cite{xiao2023streamingllm} uses attention sinks
with a sliding window---a fixed heuristic, not adaptive.

OrthoCache is unique in combining spectral eviction, single-kernel
fusion, and Lean~4 formal verification.


% ===================================================================
\section{Conclusion}\label{sec:conclusion}

OrthoCache GPU demonstrates that spectral KV-cache eviction, when
implemented as a fused Split-K Triton kernel with interleaved tile
assignment, achieves datacenter-grade attention performance on consumer
GPU hardware. The 15.3$\times$ speedup at 32K tokens and sub-linear
scaling curve represent a fundamental change in the relationship between
context length and attention latency.

\subsection{Limitations}

\begin{itemize}
\item Single-device benchmarks only (no multi-GPU NVLink validation).
\item Kernel benchmarks, not end-to-end model inference.
\item Consumer GPU (RTX 4060) --- datacenter GPU validation pending.
\end{itemize}


% ===================================================================
\bibliographystyle{IEEEtran}

\begin{thebibliography}{5}

\bibitem{orthocache_tpu}
J.~Arndt, ``OrthoCache: Hardware-Native Multi-Band Spectral Attention
Block Eviction on TPUs,'' Zenodo, 2026.
doi:~10.5281/zenodo.20518370.

\bibitem{dao2023flashattention2}
T.~Dao, ``FlashAttention-2: Faster Attention with Better Parallelism
and Work Partitioning,'' \textit{arXiv:2307.08691}, 2023.

\bibitem{kwon2023pagedattention}
W.~Kwon \textit{et~al.}, ``Efficient Memory Management for Large
Language Model Serving with PagedAttention,'' \textit{SOSP}, 2023.

\bibitem{zhang2023h2o}
Z.~Zhang \textit{et~al.}, ``H$_2$O: Heavy-Hitter Oracle for Efficient
Generative Inference of Large Language Models,'' \textit{NeurIPS}, 2023.

\bibitem{xiao2023streamingllm}
G.~Xiao \textit{et~al.}, ``Efficient Streaming Language Models with
Attention Sinks,'' \textit{arXiv:2309.17453}, 2023.

\end{thebibliography}

\end{document}
